% ===============================================================
\chapter{Quine's Two Dogmas}
% ===============================================================

\textbf{Quine}: interesting maths boi. He came about during the era of logical empiricism. There were three parts of the movement at the time:
\begin{enumerate}[label=(\alph*)]
  \item \emph{Empiricism}: generally, empiricism claims that the primary source of knowledge is sense experience. It is the primary means by which we gain (scientific) knowledge, excepting logic and mathematics.
  \item \emph{Verificationism}: sentences are meaningful only if (i) they are analytic, or else (ii) they could be confirmed or disconfirmed by experience.
  \item \emph{Anti-metaphysics}: statements which cannot be (dis)confirmed by experience are literally meaningless. These include statements of theology (`God is all-knowing'), ethics (`Giving to charity is good'), and metaphysics (`Being is the positive moment or thesis').
\end{enumerate}
Quine had some issues with verificationism specifically, as well as with the definition of analyticity.

\section{Verificationism}

\begin{definition}[Verificationism]
The meaning of a sentence consists in its method of empirical confirmation or disconfirmation. Specifically, the meaning of a sentence $S$ is its \emph{verification conditions}, i.e.\ the conditions in which $S$ would tend to be (dis)confirmed by experience. We might think of these as the set of possible experiences which would tend to confirm $S$.
\begin{enumerate}[label=\roman*.]
  \item A sentence in verificationism is \emph{analytic} if it is synonymous with a logical truth.
  \item Two sentences are \emph{synonymous} if they have the same verification conditions. This implies that analytic sentences are confirmed by all possible experiences.
\end{enumerate}
\end{definition}

\subsection*{Objections to verificationism}

\begin{objection}[Meaning is not verifiability]
We cannot know if $S$ can be disconfirmed by an experience before knowing its meaning $\Rightarrow$ the meaning of $S$ cannot be its verification conditions.
\end{objection}

\begin{objection}[Verificationism is self-refuting]
There is no possible set of experiences to confirm verificationism, at least none that we can zero in on $\Rightarrow$ either verificationism is meaningless or analytic, and there is absolutely no reason for verificationism to be analytic.
\end{objection}

\begin{objection}[Quine's objection: the fun one]
Verificationism presupposes that meaningful sentences can be divided into two kinds: those which are (dis)confirmed by experience and those which are analytic. But there are no analytic sentences. In fact, there is no such thing as meaning!
\end{objection}

\section{Synthetic vs analytic sentences}

\begin{definition}\leavevmode
\begin{enumerate}[label=\roman*.]
  \item \emph{Analytic sentences}: true in virtue of meaning alone.
  \item \emph{Synthetic sentences}: true in virtue of how the world happens to be, i.e.\ can be confirmed by experience.
\end{enumerate}
\end{definition}

\subsection*{Characterisations of analytic sentences, according to Quine}

\paragraph{Kant.}
\begin{enumerate}
  \item A subject-predicate sentence is analytic if the concept expressed by the subject is ``already contained'' within, or is subordinate to, the concept expressed by the predicate (coextensional).
  \item A sentence is analytic if its denial results in a contradiction. For example, if we deny that ``bachelors are unmarried'', then it contradicts the definition of a bachelor (``bachelors are unmarried'') and therefore results in a contradiction.
\end{enumerate}

\paragraph{Frege.}
\begin{enumerate}
  \item A sentence is analytic if it is a logical law or else is derivable from logical laws and suitable definitions alone.
\end{enumerate}

\section{Quine's attack on analyticity}

\begin{definition}[Quine's Socratic assumption]
A characterisation of analyticity is acceptable only if it takes the form
\[
  \text{A sentence $S$ is analytic iff it is \underline{\hspace{4em}}},
\]
where the blank is filled with something independently understandable, and the definition is not circular (i.e.\ it does not rely on analyticity to be understood).
\end{definition}

Quine's objection is that all three assumptions fail to meet the above definition:
\begin{enumerate}[label=(\alph*)]
  \item Kant's first characterisation, ``already contained'', is metaphorical. It applies only to sentences with ``general terms''.
  \item Kant's second characterisation, contradiction, also collapses. Suppose we take this to be true, and we want to check if the phrase ``bachelors are unmarried'' is an analytic phrase. The denial of the phrase is not directly in the form of a contradiction, but we know that due to synonymy (which we do not have any notion of), it contradicts the statement that ``bachelors are unmarried''. Therefore, we need a definition of analyticity (statements that must be true) to create a contradiction, leading to circular reasoning.
  \item Frege's characterisation collapses due to a reliance on ``suitable definitions'', which is unclear and unspecified.
\end{enumerate}

\section{Quine's attack on synonymy}

These are the characterisations of synonymy according to Quine. Some of our critiques earlier when it comes to substitutivity are criticisms Quine shares:
\begin{enumerate}
  \item \emph{Extensional substitutivity}: two expressions are synonymous if and only if they are substitutable salva veritate.
  \begin{enumerate}[label=(\alph*)]
    \item We have already proved that SSV collapses if put into an intensional language like English (there are many examples of SSV failing due to propositional attitude\ldots\ yeah, I'll save it). Quine formalises this by saying that when it comes to an extensional language, SSV guarantees only coextensiveness and not synonymy.
  \end{enumerate}
  \item \emph{Necessary substitutivity}: two expressions are synonymous if and only if they are necessarily substitutable salva veritate.
  \begin{enumerate}[label=(\alph*)]
    \item I guess we control when we can use SSV, but then this makes synonymy just as vague as analyticity.
  \end{enumerate}
\end{enumerate}

\section{Consequences of these attacks}

\paragraph{No synthetic/analytic distinction.}
\begin{enumerate}
  \item We don't know what is analytic anymore, or what is true in virtue of meaning alone. We can no longer distinguish sentences like these from statements derived from empirical fact.
  \item This undermines verificationism.
\end{enumerate}

\paragraph{Skepticism about meaning.}
\begin{enumerate}
  \item Skepticism about analyticity $\Rightarrow$ skepticism about meaning.
  \item There is no fact that governs meaning.
\end{enumerate}

\section{Confirmation holism}

Confirmation holism is one outlet of Quine's skepticism about meaning. We can hold two kinds of belief structure:

\begin{definition}\leavevmode
\begin{enumerate}
  \item \emph{Foundationalism}: there are certain primal beliefs that cannot be justified but justify other beliefs, and we build our belief system on them.
  \item \emph{Holism}: our beliefs form an interconnected web. No belief is immune to rational revision; they are justified together. However, there is a scale of how prone certain beliefs are to rational revision (those on the periphery of the interconnected web).
\end{enumerate}
\end{definition}

\paragraph{Quine's confirmation holism.}
\begin{enumerate}
  \item Quine was an empiricist, which means he believed that experience was the primary method of obtaining knowledge.
  \item Combine the above with holism, and you get a holistic belief system where rational revision happens through experience. Experience confirms the entire system of beliefs (the web).
  \item Due to the nature of the above, there is no analytic/synthetic distinction between beliefs.
\end{enumerate}

\section{Quine's meaning eliminativism}

Meaning eliminativism is the more direct outlet of Quine's skepticism about meaning.

\subsection{Radical translation}

\begin{example}[Thought experiment]\leavevmode
\begin{enumerate}[label=(\alph*)]
  \item I speak Dingleberry and walk up to you.
  \item I point at an anthill and go ``globlglob''.
  \item You will have literally no idea what I am talking about if that is the first time you've heard me speak Dingleberry. Am I talking about the ants? An anteater? An anthill? Pain? All very valid hypotheses.
  \item You call over somebody else. They point and either go ``globglobg'' or ``binglebongle''. You are still left confused.
  \item Quine repeated this thought experiment in a more interesting way, and said that any hypothesis of meaning is vastly undetermined by the evidence.
\end{enumerate}
\end{example}

\subsection{Meaning eliminativism}

\begin{definition}[Meaning eliminativism]
There is no such thing as linguistic meaning, as shown from the above, assuming our system of belief is confirmation holism.
\end{definition}

\begin{argument}\leavevmode
\begin{premises}
  \item[P1)] If there is no empirical fact as to what expressions mean, then meanings should not feature as posits of our best scientific theories.
  \item[P2)] We ought to accept the existence of only those things that are required to make our best scientific theories true, including what they posit (from Quine's ontology).
  \item[C2)] Therefore, we deny the existence of meaning.
\end{premises}
\end{argument}

\subsection{Objections to Quine's meaning eliminativism}

\begin{objection}[From logical truths]
As a logician, Quine accepts the existence of logical truths. But it is arguable that the notion of logical truth presupposes a notion of meaning. That's because there are sentences which have the form of a logical truth but contain words that are lexically ambiguous.
\end{objection}

\begin{objection}[From necessity]
Quine argues that synonymy depends on necessity, which is just as unclear as analyticity. Kripke later clarified necessity using modal logic.
\end{objection}

\begin{objection}[From modern science]
Quine believes that we must deny the existence of meaning in order to accept the empirical facts that contribute to our best scientific theories. However, linguistic semantics, a scientific field, posits meanings as a part of its best theories. Hence, a contradiction.
\end{objection}
